\scanpage{35}
les extensions résiduelles de $S'\longrightarrow S$ sont triviales. \deleted{et que les composantes irréductibles de $S'$ sont \unsure{géom.} unibranches} \inserted{Alors pour tte composante connexe $S'_i$ de $S'$, $S'_i\mathbin{-}S$ est une immersion} \deleted{\unsure{si sont}} fermée. Si $S$ est irréductible, une des immersions fermées \deleted{\unsure{est}} est un isom (car la somme des images est tout) et on gagne. \deleted{Cela \unsure{montre alors} que} \inserted{On utilise aussit un}

\deleted{pour tte composante}

\uline{Lemme} Soit $\widehat X$ schéma formel propre sur $\widehat S$ ($S=\operatorname{Spec}R$, $R$ anneau loc noeth complet), où $S$ réduit. Supposons que pour tte composante irréductible $S_i$ de $S$ (avec structure induite réduite) $\widehat X\times_{\widehat S}\widehat S_i$ soit algébrique. Alors $\widehat X$ est algébrique.

Facile par récurrence sur nb de composantes irréductibles, on est ramené au cas $S=S'\cup S''$, $S'$ et $S''$ sous-schémas fermés, et $\widehat X_{\widehat S'}$ et $\widehat X_{\widehat S''}$ \deleted{\unsure{propres}} algébriques, provenant de $X'$ et $X''$. Alors $X'_T\simeq X''_T$ ($T=S'\cap S''$) par th. de comparaison, et on recolle $X'$ et $X''$ par cet isom. On voit facilement que le $X$ recollé marche.

$(i')\Longrightarrow(i)$ Cela résulte du lemme suivant (sans doute cas particulier d'un \deleted{lemme} \unsure{résultat|vieillot} à démontrer avec \unsure{complexe|compxe} cot. relatif $\ldots$)

\uline{Lemme} Soit $\widehat X$ schéma formel \inserted{\uline{propre et}} lisse sur $\widehat S$, où $S=\operatorname{Spec}R$, $R$ anneau local noeth complet. Soit $S'=S_{\mathrm{red}}$, et supposons que $\widehat X_{\widehat S'}$ soit algébrique. Alors $\widehat X$ l'est aussi.

\mspage{34}
